\chapter{Технологический раздел}

%На листинге~\ref{lst:usbconfig.c} представлены параметры отслеживаемого 
%\includelisting
%    {usbconfig.c}
%    {Параметры отслеживаемого USB-накопителя}

\section{Выбор языка и среды программирования}

Для разработки серверной части программно-алгоритмического комплекса был выбран язык
программирования C++~\cite{Cpp}, что обусловлено следующими факторами:
\begin{enumerate}
    \item встроенные возможности объектно-ориентированного программирования, позволяющие
создавать легко расширяемые приложения;
    \item совместимость с языком C, возможность прямого вызова функций монтирования устройств;
    \item поддержка многопоточности, необходимой для параллельной обработки событий;
    \item стандартная библиотека предоставляет множество контейнеров, упрощающих разработку.
\end{enumerate}

Для разработки клиентской части был выбран язык TypeScript~\cite{TypeScript} и фреймворк Vue.js~\cite{VueJS}. 
Использование TypeScript обеспечивает статическую типизацию, что позволяет уменьшить количество ошибок на 
этапе разработки и повысить удобство сопровождения кода. Кроме того, язык предоставляет современные 
средства разработки веб-приложений и хорошо интегрируется с большинством популярных frontend-фреймворков.

Выбор Vue.js обусловлен компонентным подходом к разработке пользовательского интерфейса, что упрощает 
повторное использование элементов приложения и повышает удобство сопровождения проекта. Дополнительно 
фреймворк предоставляет встроенные средства реактивного обновления интерфейса, позволяющие эффективно 
отображать изменения состояния приложения в режиме реального времени.


В качестве среды разработки был выбран Visual Studio Code~\cite{VSCode} --- кроссплатформенная среда,
функциональность которой может быть увеличена с использованием дополнительных расширений.

\section{Выбор средства тестирования кода}

Для тестирования был выбран фреймворк GoogleTests~\cite{GoogleTests}. Данный продукт позволяет
создавать unit-тесты, e2e-тесты и интеграционные тесты. Разработан специально для C++, и использует 
стандартный его синтаксис. Позволяет реализовывать Mock-объекты для изоляции зависимостей при тестировании 
отдельных компонентов системы. Также фреймворк поддерживает параметризованные тесты, что позволяет 
проверять поведение методов на различных наборах входных данных без дублирования кода тестов. 


\section{Выбор встраиваемой системы управления реляционными базами данных (СУБД)}

Для сравнения были выбраны следующие известные системы:

\begin{enumerate}
    \item SQLite~\cite{Sqlite};
    \item DuckDB~\cite{DuckDB};
    \item H2 Database~\cite{H2Database};
    \item Apache Derby~\cite{ApacheDerby}.
\end{enumerate}

Результаты анализа существующих решений приведены в таблице~\ref{tab:compare_sudb_emb}.

\begin{table}[H]
    \caption{Результат анализа существующих решений}
    \begin{tabular}{|p{8cm}|c|c|c|c|}
    \hline
    \diagbox[width=8.5cm]{\textbf{Критерий}}{\textbf{Номер решения}} &
    \textbf{1} & \textbf{2} & \textbf{3} & \textbf{4}\\
    \hline

    Поддержка Linux
    & + & + & + & + \\ \hline

    Поддержка C++
    & + & + & - & - \\ \hline

    Бесплатное распространение
    & + & + & + & + \\ \hline

    Поддержка внешних ключей
    & + & + & + & + \\ \hline

    Поддержка ограничений
    & + & + & + & + \\ \hline

    Наличие обновлений и поддержки
    & + & + & + & - \\ \hline

    Отсутствие внешних зависимостей
    & + & + & - & - \\ \hline

    Размер библиотеки(Мегабайты)
    & 1.5 & 70 & 2.5 & 6.1 \\ \hline

    \end{tabular}
    \label{tab:compare_sudb_emb}
\end{table}

На основе анализа, для разрабатываемого комплекса больше всего подходит SQLite.